% ============================================================
%  Georg Brandes, Dualismen i vor nyeste Philosophie
%  Dualism in Our Most Recent Philosophy
%  Kjøbenhavn: Forlagt af den Gyldendalske Boghandel (F. Hegel), 1866.
%  76 pp., antiqua.
%
%  TRANSLATION (English) --- Hans Halvorson
%  Translated from transcription.tex (Danish), which is COMPLETE and
%  image-verified against two independent scans of two different copies
%  (measured error rate 0.00/page over a seeded random fifth of the body).
%  Method: ../../../TRANSLATION-PLAYBOOK.md.
%
%  ---- WHAT THIS BOOK IS ----
%  Brandes's first book, published at twenty-four: a polemic against Rasmus
%  Nielsen's doctrine that faith and knowledge are ABSOLUTELY HETEROGENEOUS,
%  i.e. principles of such different kinds that they cannot contradict one
%  another and may be held together in one consciousness without discord.
%  The argument is that the doctrine cannot be held: if faith and knowledge
%  are utterances of one and the same being, they cannot be utterances of
%  absolutely heterogeneous principles, and the heterogeneity that is
%  supposed to keep them from colliding is exactly what would make a single
%  consciousness impossible. What the doctrine produces is not peace but a
%  division of the person --- one man in his science, another in his life.
%
%  ---- REGISTER ----
%  Readable modern English. Brandes at twenty-four writes long periodic
%  sentences in the Danish manner of the 1860s; the periods are broken where
%  English will not carry them, and the idiom is modernised, but nothing is
%  softened: the book is sarcastic, and the sarcasm is the argument's manner,
%  not decoration. Where the Danish is ironic the English is ironic.
%
%  ---- TERMINOLOGY ----
%  Aligned with the Nielsen and Brøchner translations in this repo, since the
%  book is an intervention in their controversy and the three must read
%  together. Terms flagged once here rather than glossed on the page:
%
%    Uensartethed        = heterogeneity   (absolut Uensartethed = absolute
%                          heterogeneity --- the doctrine under attack; it is
%                          Nielsen's own term and Brøchner's, and the English
%                          is the same in all three files)
%    Viden               = knowledge
%    Erkjendelse         = cognition       (as in nielsen/propaedeutiske-logik;
%                          NOT "knowledge", which is reserved for Viden. The
%                          two words are both frequent here and the argument
%                          turns on holding them apart)
%    Videnskab           = science         (Wissenschaft: organised knowledge
%                          of any kind, incl. the humanities --- never
%                          "natural science", for which Brandes writes
%                          Naturvidenskab)
%    Fagvidenskaberne    = the special sciences
%    Bibeltroen          = biblical faith
%    Tro / Troen         = faith
%    Bevidsthed          = consciousness
%    Forestilling        = representation
%    det Antirationelle  = the anti-rational
%    det Suprarationale  = the suprarational
%    Under / Underet     = miracle / the miracle
%    Speculation         = speculation (the Hegelian kind; pejorative here)
%    den Dannede         = the educated man; Dannelse = cultivation
%    Livsanskuelse       = view of life;  Verdensanskuelse = world view
%    Personlighed        = personality
%    Grundideernes Logik, Philosophisk Propædeutik, Indledning til Ethiken:
%                          work-titles stay in Danish, as the author cites them
%
%  ---- CONVENTIONS ----
%  - Quotes. Danish „...“ -> English curly doubles ``...''.
%  - Emphasis. \emph{} in the Danish is letterspacing in the print (the book's
%    ONLY emphasis device --- there is no italic in it anywhere); every one is
%    carried over and renders italic here.
%  - Latin and foreign phrases -> \textit{} (principiis obstare, suum cuique).
%  - Footnotes. 22 in the book, marked *) and **) as printed; the \fnmark
%    macro of the transcription is reproduced so the marks match page by page.
%  - Page markers \opage{N} and "% --- p. N ---" carried at the same joints,
%    so the two files can be read side by side.
%  - The old „ɔ:“ id-est mark -> "i.e."
%
%  ---- DEFECTS: DANISH AS PRINTED, ENGLISH CORRECT ----
%  The transcription is diplomatic and carries every printer's error; the
%  translation reproduces the SENSE and logs the divergence in a % note at the
%  site. Sites: p. 9 `bos' for `hos' (read "with"); p. 18 `teleolologisk' for
%  `teleologisk' (the doubled syllable cannot be reproduced in English and is
%  translated as the correct word); pp. 20 and 24 unpaired closing quotes (the
%  English balances at 0).
% ============================================================
\documentclass[12pt,a4paper]{book}
\usepackage[utf8]{inputenc}
\usepackage[T1]{fontenc}
\usepackage{libertinus}
\usepackage{libertinust1math}
\usepackage{textalpha}
\usepackage{graphicx}
\usepackage[english]{babel}
\usepackage[protrusion=true,expansion=true]{microtype}
\usepackage{setspace}
\usepackage[margin=1.2in]{geometry}
\usepackage{enumitem}
\usepackage{fancyhdr}
\setlength{\headheight}{28pt}
\addtolength{\topmargin}{-13.5pt}

% Footnotes as the 1866 printing sets them: *) and **), not 1, 2, 3.
% Same mechanism as transcription.tex: the mark is set explicitly at each
% site with \fnmark{*)}{...} or \fnmark{**)}{...}, which restores the default
% afterwards so a missed one cannot silently propagate.
\renewcommand{\thefootnote}{*)}
\newcommand{\fnmark}[2]{\renewcommand{\thefootnote}{#1}\footnote{#2}%
  \renewcommand{\thefootnote}{*)}}

\usepackage{hyperref}
\hypersetup{
  pdftitle={Dualism in Our Most Recent Philosophy},
  pdfauthor={Georg Brandes},
  hidelinks
}

\pagestyle{fancy}
\fancyhf{}
\fancyhead[LE,RO]{\thepage}
\fancyhead[RE]{\textit{Dualism in Our Most Recent Philosophy}}
\fancyhead[LO]{\textit{\leftmark}}

\setstretch{1.3}

\title{\textbf{Dualism}\\[6pt]
  \large in Our Most Recent Philosophy}
\author{G.\ Brandes}
\date{Copenhagen: Published by the Gyldendal Bookshop (F.\ Hegel), 1866\\[10pt]
  \small Translated from the Danish by Hans Halvorson}

% ---- original-page markers ----
% \opage{N} marks the point at which printed page N begins, at the same word
% as in transcription.tex.  Because English runs at a different length the
% marker falls mid-sentence at the same JOINT, not at the same position on
% the typeset page.
\usepackage{marginnote}
\newcommand{\opage}[1]{\marginnote{\footnotesize\textit{[#1]}}}
\newcommand{\apage}[1]{\marginnote{\footnotesize\textit{[#1]}}}
\begin{document}
\maketitle
\thispagestyle{empty}
\newpage

\tableofcontents
\newpage

%% ============================================================
%%  PRELIMINARIES (printed pp. 1--6)
%% ============================================================
\chapter*{Title Page and Contents}
\addcontentsline{toc}{chapter}{Title Page and Contents}
\label{ch:prelim}

% --- p. 1 ---
% Half-title.  Both lines letterspaced in the print, `Dualismen' larger.
\begin{center}
\opage{1}\emph{\Large Dualism}\\[10pt]
\emph{in Our Most Recent Philosophy.}\\[14pt]
\rule{2cm}{0.4pt}
\end{center}

% --- p. 2 ---
% Blank in the print.
\opage{2}

% --- p. 3 ---
% Title page.  The imprint is left in Danish below the translated display
% lines: it is the book's own imprint, not matter to be Englished.
\begin{center}
\opage{3}{\Large\textbf{Dualism}}\\[10pt]
\emph{in Our Most Recent Philosophy.}\\[14pt]
\textbf{By}\\[10pt]
\emph{G.\ Brandes.}\\[16pt]
[Publisher's device]\\[16pt]
\emph{\textbf{Kjøbenhavn.}}\\[4pt]
Forlagt af den Gyldendalske Boghandel (F.\ Hegel).\\[4pt]
{\small\emph{I.\ Cohens Bogtrykkeri.}}\\[6pt]
1866.
\end{center}

% --- p. 4 ---
% Blank in the print.
\opage{4}

% --- p. 5 ---
% Indhold.  The entries carry no chapter numbers and end in no full stop,
% though the chapter heads in the body do both; reproduced as in the Danish.
\begin{center}
\opage{5}\emph{\textbf{Contents.}}\\[4pt]
\rule{1.5cm}{0.4pt}
\end{center}

\begin{flushleft}
\hfill\textbf{Page}\\[2pt]
Introduction \dotfill 7\\
Absolute Heterogeneity \dotfill 25\\
The Anti-Rational \dotfill 47\\
The Consequences of Dualism \dotfill 55
\end{flushleft}

\begin{center}
\rule{1.5cm}{0.4pt}
\end{center}

% --- p. 6 ---
% Blank in the print.
\opage{6}

%% ============================================================
%%  1. INTRODUCTION  (printed pp. 7--24)
%% ============================================================
\chapter*{1.\ Introduction.}
\addcontentsline{toc}{chapter}{1. Introduction}
\markboth{Introduction}{Introduction}
\label{ch:1}

% --- p. 7 ---
\opage{7}There is something in Prof.\ R.\ Nielsen's account of the relation
between science and biblical faith that is bound to win him many hearts. One
feels a reassurance joined to satisfaction, a glad confirmation of something
one had wished for, perhaps without much hope, on hearing a thinker speak so
severely of the impotence of human thought, on seeing so great and
comprehensive a philosophical talent so conscious of philosophy's limits. One
came to the philosopher all too accustomed to the phrase, repeated to the point
of triviality, that the educated man --- sometimes with pride, sometimes with
pain --- feels himself torn loose from the presuppositions on which the
unlearned and simple build with confidence; and one goes away strengthened in
one's trust that it is only superficial and immature philosophy that leads away
from biblical faith, and not the true, deep, real understanding of what can be
understood. And admiration mixes itself into the satisfaction. This is the same
man who with so much zeal and warmth has upheld the independence of the special
% --- p. 8 ---
\opage{8}sciences against the pretensions of an arrogant speculation; who,
standing at the high point of his age's scientific culture, girded with the
strength-belt of his genius and his cultivation, has moderation, sobriety,
humility and piety enough to declare that life is more than contemplation,
character more than intelligence, practice more than theory, and religion
higher than science. Other thinkers, lesser thinkers, foreign philosophers and
German ones above all, have let themselves be lured into the ill-famed modern
freethinking; they have given offence by attacking what the tradition of
centuries has stamped as holy. Our own philosophy, under this thinker's
auspices, will run in a better track: without renouncing the freedom that is
thought's right and its condition of life, it will preserve reverence for the
holy in a dignified self-command, clean of the excesses to which science has
given itself over everywhere else in Europe and in part here in Denmark too.
Who has not had enough of barren polemics, of the appetite for demolition and
the rage for annihilation? What is finer and better than this so sincerely and
vigorously declared effort to make peace at last between powers that, each of
infinite importance in itself, ought to be working peaceably together toward a
common goal!

In these or similar terms admiration will clothe itself, and the circle that
talks this way is hardly a small one. For this philosophy is in more than one
respect a philosophy of reconciliation. It reconciles
% --- p. 9 ---
\opage{9}not only principles but parties; it looks for its friends in opposite
camps. Though in religious matters its whole content of ideas comes from
Kierkegaard, it is nonetheless raised above what it must regard as a limitation
% The Danish reads `bos' for `hos' --- a printer's error, transcribed as
% printed there; translated here as the correct word ("in", of that genius).
in that great genius. It sees not merely, as nearly everyone does, what is in
many ways stirring and good in the life that has issued from Grundtvig and his
followers; it has an eye for the deep truth in Grundtvig's discovery, the
standpoint for which Kierkegaard, and so many with him, are unable to find. For
this philosophy knows how to value fresh immediacy as well as fully
carried-through reflection. It is hostile only to the sphere in between, only
to the murky commingling of the scientific and the religious --- in one word,
to theology. And yet it is convinced that the theologians, once they get a
proper view of what is in the highest sense their own interest, will feel the
need to gather under its broad wings together with their former opponents; for
what the best among these want, it wants too: the holiness of the uncut,
inviolable word of the Bible, and faith's independence of an ever more
demanding and inconsiderate modern science. Finally, it does not shut itself up
with its initiates in a sealed room. It is intent on gaining currency among the
generally educated, and indeed among the great teachable mass of the people,
whom it is a pleasure to instruct for anyone lucky enough to have the gift for
it. It is as popular
% --- p. 10 ---
\opage{10}when it turns outward as it is strict and exact in its dealings
within; as soft and conciliatory when the business is to unite whatever is
great and distinguished that our fatherland has brought forth in godly and
churchly matters\fnmark{*)}{See e.g.\ Nord.\ Universitetstidskr., supplementary
number to vol.\ 10, 1866, p.\ 75.}, as it is firm of character and inexorable
toward theology's attempts at mediation (which are a matter of indifference to
the general public).

How much there is in such a programme that appeals to what is best in a man,
and that would at any time make an impression on every mature and developed
mind --- and how much there is in it, finally, that must strike precisely the
people of our time and our country!

The nineteenth century turned early against the philosophy of the eighteenth;
but here in Denmark perhaps more decidedly, and at all events more
persistently, than anywhere else. However fiercely our theologians have warred
with one another, Danish literature bears neither many nor strong traces of the
critical movement that characterises German and French science in the present
century. The age of the rebirth of religious movements has nothing to spare for
rationalism. It abhors the flatness of the old rationalism; it strongly doubts
that any good can come of unsettling the authority of positive religion once
more; it is reluctant to see any of revelation's outworks given up, for it
feels the necessity of \textit{principiis obstare}. This
% --- p. 11 ---
\opage{11}turn of mind Prof.\ Nielsen comes out to meet. He secures the
outermost entrenchment of the positive by setting a yawning abyss between
science and religion, and throwing the theologians headlong into it.

But the nineteenth century is also the age of natural science and of the great
discoveries. If it is clear about anything, it is that in the domain of science
we must have the field free, we must be able to investigate and determine
unchallenged by the preconceived and narrow judgments of an outdated orthodoxy.
This consciousness too Prof.\ Nielsen comes out to meet. When the talk is of
science, then the talk is not of faith; then the Professor stands --- or at any
rate means to stand --- free and unhindered by any binding prejudice whatever,
and no freethinker, however impatient, has been able to speak more vehemently
on the point than he.

What more can one want, what can one still wish for beyond this magnificent
execution of \textit{suum cuique}! Who will raise an objection here --- unless
perhaps the straggler who has not yet grasped that the age of dogmatics is
over, or the narrow-minded man who cannot bear that the great public should be
given a share in philosophy's results, or the individual who comes forward not
for the sake of the cause but to push his own little self into the foreground
with the help of the great cause\fnmark{*)}{Cf.\ R.\ Schmidt: Om
Selvmodsigelsen i Prof.\ Nielsens Lære, p.\ 44.\endgraf
Just as Mr.\ Schmidt supposes that something morally incriminating must lie
behind the act of standing up against Prof.\ Nielsen, so Mr.\ Henrik Scharling
cannot conceive that one might speak against theology on any but personal
grounds. He declares that I appear to be of the opinion that if theology cannot
make me a Christian, then it is barren and impractical; and he adds
instructively that one can of course no more become a Christian merely by
studying theology than one can recover one's health merely by studying
medicine. Although Mr.\ Scharling's drollery as literature's mentor is not
unknown, and although one must no doubt be prepared to meet a certain
familiarity with omniscience in a theological weekly, this outburst of Mr.\
Scharling's was nevertheless a welcome surprise to me. It is so agreeable to
find in a journal information about one's own spiritual history --- and
information, at that, of which one was previously quite unaware.}?

% --- p. 12 ---
\opage{12}Let us turn the medal over for once and look at the other side. This
doctrine, which appeals to so much that is good in people and shows itself so
much in tune with the times --- might it not also be addressing itself to
something less good, to what is lower in us, indeed to the lowest? We human
beings are by nature inclined to slackness, to dividedness, to dissipation. Most
people are happy to have one view of life in joy and another in sorrow, one at
home and another in church. The teacher prefers to profess one principle when
he speaks and another when he acts; the priest has trouble living as he
preaches; the merchant generally has one morality in business and another in
life; and the poet sings again and again of
% --- p. 13 ---
\opage{13}ideals that are least of all to be met with in his own conduct. The
one sphere, after all, is no concern of the other; the one hand need not know
what the other is doing. And now the man of science! Brought up in some
definite church faith, influenced from childhood by some definite accidental
circle of representations --- the one that prevails in his church, in his
country and in the town where he lives --- he encounters in science facts that
conflict with this faith. Here is a battle to be fought out, a victory to be
won. But could peace not be had more cheaply? Appalling thought, that the faith
preached in Ebeltoft or in Paris or in St.\ Petersburg should not be the
highest truth! Appalling faithlessness, to break with a faith in which one's
Ebeltoft or Parisian or St.\ Petersburg parents placed all their trust! What is
more natural, then, than the practical expedient of being one man in science
and another in one's life --- of living one's whole life through in a chaotic
world view with fragments of the past's outlook and scraps of the present's,
while appearing in one's writings dressed in the scientific drapery of the
nineteenth century.

Yet however naturally this comes, there is something painful in such halfness.
One is reluctant, no doubt, to gather oneself for the forceful, single
decision; but neither does one care to waver without principle between opposed
points of view. No compromise! --- that is the watchword nowadays; far rather a
sturdy disharmony! Halved one will be, if it cannot be otherwise, but then
% --- p. 14 ---
\opage{14}one will be halved thoroughly. What a relief, then, for the hard-pressed
soul is the new doctrine, which does grant science its right, if under
compulsion, but all the same leaves me my childhood faith --- and my nursery
learning too, if it comes to that --- in good order; which in short raises
unprincipledness itself into a principle, and gives the uncertain man precisely
what he craves: leave to be whole in his halfness, to breathe fully as a double
human being. His understanding needs air, and it is given free play, a field to
range in without limits; while his feeling, which is at odds with his
understanding, and his will, which in its need of outward support parts only
with difficulty from any article of faith handed down from the fathers, need
give up nothing at all. Faith is justified the moment it is merely firm; its
kingdom is without barriers, its territory has no boundary in common with the
land of cognition. Each of the two realms has its own laws, its own ruler,
indeed its own God; they cannot fall out with one another over what is true,
for this very discord between them is
the truth, is the peace. As when two great powers stand fully armed against one
another, but mediation intervenes, no war breaks out, and a settlement is
struck in which each party, fully armed, reserves all its rights --- so here
the debate is stifled under a storm-laden pause: neither of the two opponents
gives up a jot. Regarded as a temporary situation, or as a truce after old
battles, this standstill may have its use; it is even of value to
% --- p. 15 ---
\opage{15}science, if science is really allowed to live its own life
undisturbed by outside interference. But if it is to be taken for a settlement
of the dispute, for the long-awaited solution of a great problem --- and this
is what it is given out to be\fnmark{*)}{See e.g.\ ``Fædrelandet'' 1866 no.\
172.} --- then its weakness is so glaring that it would be a disgrace to our
science if it were not pointed out.

We have seen what might tempt even the man of science to attach himself to the
new doctrine and to realise it for his own part. But if the temptation can
become too strong even for him, how much greater must it be for the layman.
Suppose he has from childhood a need for unconditional surrender and heartfelt
faith. Suppose, next, that while his heart seeks positive religion, his reason
is at the same time repelled by many a crudity that seems inseparably bound up
with religion. If he now rejects religion unconditionally, his heart languishes
forever in its want; if instead he accepts it without reservation in its
handed-down shape, it is as though an iron band were clapped about his brow. So
he sits fast in the dilemma of choice. Suppose, then, that one evening he steps
into a lecture hall where listeners of both sexes, educated and uneducated,
have crowded together. The lecture he happens to hear is aimed at speculative
theology. With all the assurance that learning, logical virtuosity and lively
eloquence confer, the speaker sets out
% --- p. 16 ---
\opage{16}how dangerous it is, in matters of faith, to give science so much as
one's little finger. On science's behalf he enters a protest against the
credibility of the biblical miracles. He presents first, perhaps --- always
from science's standpoint, of course --- the improbability, the absurdity of
some of the miracles that have little nimbus for his hearers' imagination but
are all the more an offence to their understanding. The listening circle is not
disquieted yet; it feels rather a relief of mind at letting its pent-up
indignation run free for once. But the speaker goes further. He approaches,
perhaps, miracles that stand to everyone as holy mysteries; for a moment it
looks as though he meant to lay hands on them, and a tremor passes through the
assembly. The consequence of the theologians' attempts at explanation has now
been made plain. Then the speaker turns the hourglass of his lecture. With warm
enthusiasm, with still greater eloquence, he presents the miracle seen in the
light of faith. And now, when the critical inclination has been allowed to
spend its rage, now that one has learned to shudder at the train of thought one
was lately willing to enter upon, one embraces with the same passion great and
small, the most exalted religious commandment and the most trifling miracle;
everything seems equally holy and precious now that the danger of losing it is
hardly past. What has the hearer not felt and lived through in such an hour,
how freely have all his faculties ranged! Reason and the craving for faith are
equally
% --- p. 17 ---
\opage{17}satisfied --- the one because it has been given leave to throw the
whole flock of miracles out of science's front gate, the other because it has
been given complete freedom to let that same throng in at the back gate that
stands open to it in the individual's practical existence. Forward and back
seem equally far; but the hearer has exercised his spiritual powers, and he has
learned once and for all how he is to handle his understanding every time its
impatience threatens to expose him afresh to the temptations of offence. In his
joy it escapes him that what has been carried out here is only a feigned
manoeuvre; that since the concept of God and the concept of the world which
Prof.\ N.\ sets against theology's are, on the Professor's own view, not the
true ones, not the ones that have reality, the battle fought here against
speculative theology is only a sham battle --- a mock skirmish, as Holberg
calls it --- and that all that has happened is a restationing of the
suprarational from the domain of science to the domain of life: a tactical
change, analogous to the one that would have taken place in chemistry if the
chemists had at some point agreed to tolerate phlogiston no longer in their
science, but only so as to be all the more fervently convinced of its existence
outside science.

Or perhaps the groping, unsettled man turns to literature, and stops one fine
day at the reading of a page like this\fnmark{*)}{R.\ Nielsen: Forelæsninger
over Philosoph.\ Propædeutik. Efteraarshalvaar 1865 p.\ 159.}, containing the
close
% --- p. 18 ---
% The Danish has the printer's doubled syllable `teleolologisk' at the first
% occurrence, letterspaced, transcribed as printed; it cannot be reproduced in
% English and is translated as the correct word, emphasis kept.
\opage{18}of a conversation in which theology's defender comes off pitifully.
The theologian C.\ has put forward the proposition that in the objective sense
everything, or else nothing, is to be called a miracle. To this A.\ replies as
follows:

``That a man walks on the sea without sinking --- is that then just as much in
accord with the immutability of the laws of nature as that someone falls into
the water and drowns? C.\ The human being, who is otherwise subject to the laws
of gravity, may very well walk on the sea when it happens
\emph{teleologically}, in virtue of a continued creation! A.\ But then, so far
as the immutability of the laws of nature is concerned, there is presumably
nothing offensive to deeper thinking either, if we consider it properly, in a
pope, Leo XII, having been able to beatify the Minorite Julianus, who in our
unbelieving, revolutionary nineteenth century made roasted birds fly? C.\ That
I for my part take the miracle of the roasted birds to be a stupid fable, a
priest's lie, comes not at all from my seeing in the event itself any breach of
the immutability of the laws of nature, but simply from my not being a Catholic
and consequently not being able to convince myself that the event happened
teleologically, in virtue of a continued creation. Were I a Catholic and had a
Catholic's faith, I should no doubt be capable of reconciling my faith with my
knowledge, my concept of creation with my concept of nature . . .''

What has happened here? In making theology's attempts at explanation
ridiculous, in showing the supposed understanding of the miracle to be merely
supposed,
% --- p. 19 ---
\opage{19}the mockery strikes speculative dogmatics not only as that which ---
in principle --- is capable of comprehending so insane a miracle, but as that
which is not in a position to turn it away. And so a ridiculous light falls
over the miraculous world itself with which dogmatics deals: the world, namely,
in which roasted birds can be made to fly. Here again the sympathetic reader
feels his understanding satisfied. Although no miracle venerable by its age has
been brought even under the shadow of a critique by the understanding, he takes
comfort all the same in the thought that he at least is not obliged to resort
to arguments like C.'s in order to fend off miracles of the sort of the one
with the birds. But this feeling rests on pure illusion. For take up the thread
of the conversation and turn the reply against Mr.\ A., and it will appear that
his arguments are not a hair better than his humiliated opponent's, and that
his position is quite as peculiar as Mr.\ C.'s. Ask him, with only a few words
altered: Is there then, so far as the immutability of the laws of nature is
concerned, anything offensive to \emph{your} innermost consciousness in a man
of our century having made roasted birds fly? Then the answer would have to run
something like this: That I for my part take the miracle to be a stupid fable,
a priest's lie, comes not at all from my seeing in the event a breach of the
immutability of the laws of nature --- for even if it does contain one, I dare
not
% --- p. 20 ---
% The Danish closes this reported speech with an unpaired “ (the opening „ is
% not on p. 19 in the print); the English balances.
\opage{20}deny its possibility --- but simply from my not being a Catholic and
consequently not believing that miracles have happened on earth since the first
century. That is why I used the words: ``in our unbelieving, revolutionary
nineteenth century''. Were I a Catholic and had a Catholic's faith, I should no
doubt be capable of reconciling my faith with my knowledge; for I am not
answerable to my reason for what I accept in faith. The immutability of the
laws of nature is a norm that holds firm only in a constitution of nature to be
found in science and on paper, and not in the actual constitution of the world,
to which I in my existence stand as an obedient subject. Against that
constitution no miracle offends which faith declares holy and which strengthens
the conviction that the whole is absolutely dependent on unconditioned
omnipotence.'' In other words: Mr.\ A.\ too is of the opinion that in and of
itself the immutability of the laws of nature is no obstacle to roasted pigeons
flying into our mouths. He does not expect it to happen, but only because the
time in which such things occurred is over.

When one points out to one of our Neo-Orthodox --- who will not grant reason or
science the feeblest right to have a say about religion's object --- that in our
day it will not do to hold fast to stories such as the one about the origin of
languages from the deity's wrath at a presumptuous tower-building, or about the
rainbow's coming into being as a bodily thing on the occasion of a definite
historical
% --- p. 21 ---
\opage{21}or unhistorical event, except as beautiful and poetic legends, the
person in question generally answers that the Bible has no intention of
teaching us natural science or philology, that it has no intention of enriching
our knowledge at all, that this is a wretched misunderstanding, and so on. This
answer plainly rests on a mistake. What is said here of the Bible does not hold
of the Bible, but of religious proclamation, of the ethical commandment. That
does not address itself to our desire for knowledge, and has no intention of
imparting to us any information about the origin of languages or of natural
phenomena. What the Bible \emph{intends} to teach us is hard to
determine\fnmark{*)}{That such an expression is frequently heard is only one
among many testimonies to the highly incorrect usage that has crept in among
us, by virtue of which people speak, among other things, of \emph{faith,
religion, science, the will} as though these were personal beings. Examples of
this metaphor-philosophy are furnished in the richest measure by Dr.\
Heegaard's ``Indledning til Ethiken''. Just one here: when on p.\ 446 he
declares that the will's concentration out beyond the rational is a step in
which the will raises for itself no theoretical scruple whatever, one is put in
mind of the old saying: virtue is no peach. One might with the same right speak
of steps in which the heart does not trouble its head in the least.}. But ask
simply whether the Bible does not --- whether it intends to or not --- teach us
something of the kind, and the answer is yes; from which it follows that the
orthodox man, if he gives up the Bible's authority on these points, has broken
once and for all with his principle.

It might now seem as though a disciple of
% --- p. 22 ---
\opage{22}Prof.\ Nielsen would in this case be able to escape entirely the
difficulties in which the orthodox man finds himself entangled. For no one has
drawn a sharper distinction than theology's tireless antagonist between what
concerns the religious consciousness and what does not. One would expect that
these facts, being without significance for the religious life, must on Prof.\
N.'s doctrine fall to science's criticism and so lose their exceptional
standing. But no: it turns out that they are regarded as holy matters of fact.
They are preserved in full inviolability, mummies of a faith that will not give
up even what is dead; and it is --- however unintelligible it sounds --- the
ethical in man, the ethical demand, that secures these corpses against critical
dissolution. The fundamental ethical demand necessitates an ethical ideal; the
good cannot exist without one who is good, the absolutely good one only as
omnipotent, the omnipotent one only as absolutely creating, and creation in
turn can only have taken place in the six days (see Heegaard, op.\ cit., p.\
392). With this the first links are forged of the chain that must presumably be
thought to bind together all the holy facts of the Old and New Testaments. The
chain has not been completed yet, but the demonstration of the whole nexus will
not, one hopes, keep us waiting too long; and then the complete assemblage will
lie before us as a fully executed system --- which it will not be advisable for
anyone to venture to call theological. So if one asks Prof.
% --- p. 23 ---
\opage{23}Nielsen's disciple: Can one reject the Old Testament's story of the
Babylonian confusion of tongues, or the New Testament's story of the casting
out of demons, and still stand in relation to the ethical and religious ideal?
--- then in the bewilderment of the first moment he will perhaps answer yes;
but as soon as he has collected himself he will consistently deny it: the
ethical cannot do without the religious, the religious cannot do without
revelation, revelation cannot do without these presuppositions. Suppose one
then objected: Do try for once to look at the thing freely and squarely; judge
without any side considerations --- do you accept these facts or not? The
answer would necessarily run about as follows: ``A plain yes or no I cannot
give to your question; for I should otherwise be making knowledge the judge of
things that lie outside its competence. As a truthful man I assume, in the
interest of truth, that these so-called facts are quite unhistorical; but as a
moral individuality I assume, in the name of the good, that these same facts
are completely reliable and real.''

``But,'' the questioner bursts out, ``do you not see that there is a third and
simpler way of regarding such accounts? I am not asking you whether they are
true, for with the word truth all manner of tricks can be played; answer only
whether you hold that what they report ever happened. Do you believe that the
rainbow was set in the heavens, or do you not believe it? Do you believe that
those madmen were possessed by demons, or do you not?'' Then the answer comes
again:
% --- p. 24 ---
% The Danish closes this reported speech with an unpaired “; the English
% balances.  Six spaced points in the print.
\opage{24}``In the interest of the true I assume, and so forth; in the name of
the good I hold fast, on the contrary . . . . . .'' --- and with that the
conversation must be at an end. But here even the most conciliatory mind comes
to recall the little sharp verse:

\begin{verse}
``The sun of Spirit, or the parson's taper?\\
Yes or no! and none of your vapour!''
\end{verse}

\begin{center}
% The chapter-end ornament: a short centred double rule, as in the print.
\rule{1.7cm}{0.4pt}\\[1.5pt]
\rule{1.7cm}{0.4pt}
\end{center}

%% ============================================================
%%  2. ABSOLUTE HETEROGENEITY  (printed pp. 25--46)
%% ============================================================
\chapter*{2.\ Absolute Heterogeneity.}
\addcontentsline{toc}{chapter}{2. Absolute Heterogeneity}
\markboth{Absolute Heterogeneity}{Absolute Heterogeneity}
\label{ch:2}

% --- p. 25 ---
\opage{25}It falls out fortunately or unfortunately, according to the
standpoint one occupies, that Prof.\ R.\ Nielsen places at the head of his
doctrine of the relation between the religious and the scientific a
fundamental proposition on which the whole edifice rests. One does not, then,
have to hunt laboriously for the theory's governing thought; it is expressed so
definitely that it cannot possibly be missed. And if one cannot go along with
it, one need not get involved with the rest: for either that thought is carried
through consistently, in which case it will permeate the whole development and
at every point make it impossible for one to adopt it, or else it is not
unfolded consistently, and then one can wish even less for a closer
acquaintance with the details. The proposition, which in a number of articles
in ``Fædrelandet'' has been called a plain self-contradiction, runs as follows:
the same consciousness can without contradiction unite faith and knowledge as
absolutely heterogeneous principles. ``Without contradiction'' --- that is the
safety valve fitted to keep the proposition from blowing
% --- p. 26 ---
\opage{26}up, and the investigation to be undertaken here aims only at testing
whether the ventilation is sufficient. The proposition put forward is a
fundamental one; but it is hardly too bold to ask what it rests on. It must
presumably rest on the general propositions \emph{that two absolutely
heterogeneous utterances can in general be utterances of the same thing} ---
faith and knowledge are, after all, utterances in principle, i.e.\ utterances
of one being --- \emph{and that two absolutely heterogeneous principles can in
general} --- shall we pass over ``without contradiction''? --- \emph{be united
in one and the same consciousness}. Suppose for a moment that this last
proposition were true: why then does it say ``can'' and not ``must''? If they
can be united, then they can also fail to be united, and so there is no
necessity that compels. And does not the same hold when one takes the two
definite principles, faith and knowledge? Beyond doubt: one can unite them and
one can leave it be, i.e.\ the union has no universal validity. But if that is
how matters stand, what brings it about that the man who does unite them
really unites what can be united? The question here is whether that can be made
\emph{clear}. Since the quarrel and the reconciliation between faith and
knowledge have been carried by Prof.\ Nielsen out of the domain of science and
over into that of ``existence'', it has become hard to keep an eye on how the
thing is actually conceived; for of existence one may use the old phrase that
Hegel used of Schelling's absolute: it is the night in which all cats are grey.
% --- p. 27 ---
\opage{27}The so-called existential, which lays so strong a claim to count as
life's true full-blooded concretion, is in reality just the opposite: a pure
abstraction, namely what is left over when one subtracts from the existence of
the human spirit \textit{in toto} everything theoretical. It is therefore just
as far from being true that the human being who exists \textit{sensu eminenti},
in the current philosophical jargon, leads a true and genuine human life, as it
is that a dog from which a naturalist has experimentally removed part of the
brain can be said to live its life wholly and really. What goes on in
``existence'' can accordingly be checked neither easily nor with certainty. But
since Prof.\ Nielsen wants clarity above all, it will of course not do for him
at any point --- and least of all at the decisive one --- to put out the light
so that nobody can see what is happening, and to light it again only once
whatever was to be carried in and out in the theory's interest has been carried
in and out. What is it, then, that brings the possible union into being? Why,
what is there to resort to here but the crassest immediacy, personal need,
which at least \emph{can} be regarded as an entirely individual matter. So this
is the philosophical marvel that so large a chorus has hailed: a higher
brutalism on the basis of instinct, of immediate assurance, or something else
of the sort. If we are to accept this doctrine, let us first give up being
spirits!

In taking refuge in immediate
% --- p. 28 ---
\opage{28}practice, which as such is without universal validity, one is
entitled to satisfy one's personal need as one pleases. But although the
consequence is undeniably this, Prof.\ Nielsen does not draw it outright. What,
then, is the source from which the philosopher draws the content that fills out
one of the two empty frames set up in the general proposition? It is
revelation, scripture, religion --- by which is meant, without further ado, the
Christian revelation, scripture and religion as these are understood in
Protestant countries, and as they have latterly been understood in Denmark in
particular. The man of science who takes Prof.\ Nielsen's theory for his guide
can manage this only by conducting himself as a subtle spirit and then, in a
twinkling, turning into a charcoal-burner who answers one's objections with a
broad smile: what a man believes, he believes, and no talk about it --- and
then, when the gust has blown over, throwing off the jerkin, washing, combing
his hair, putting on a black coat and being a subtle spirit again. If anyone
should commit the blasphemy of asking what \emph{concept} one is to form of a
revelation in general, and of the Christian one in particular, then one is a
charcoal-burner and says: what stands written, stands written. If it says that
one is to bake one's own bread, then one is to bake one's own bread; if it says
nothing about brewing one's own beer, then one need have no scruples about the
beer. If it says that one is to repent one's
% --- p. 29 ---
\opage{29}sins, then one is to repent one's sins; if it says that God created
the world in 6 days, then he did it in 6 days, for had it said that he had done
it in 7 or 14 days, then he would have done it in 7 or 14 days, and so on. (See
Dr.\ Heegaard's presentation of Prof.\ Nielsen's doctrine, ``Indledn.\ til
Ethiken'' p.\ 392.) Revelation is a paradox, an absolute paradox, and as such
without any rationality whatever. So do not ask: what is a revelation? Or: may
one draw consequences from what is given in a revelation? Or: how far may one
draw them? Such questions are asked by reasonable people, but there is no art
in being reasonable. The art is, the moment such questions are raised, to clap
the wishing-cap on one's head --- and then one is a charcoal-burner, if that is
what one wants. For it goes without saying: if one is entitled to go so much as
a single line beyond the boundaries drawn by revelation's words (6, not 7 or
14), then the extraordinary vanishes and the standpoint has been given up. ---
So: one is a charcoal-burner, and that is the one side of the matter. If one
stopped there, one might still be an object of sympathy and respect. Respect
for the charcoal-burner who is in truth what he calls himself! Honour to him
when, like the charcoal-burners in Ploug's ``Fredericiaslaget'', he goes
straight at his mark and strikes for his banner with charcoal-burner's fists
and charcoal-burner's faith! But when, at the same time as he is a
charcoal-burner, he is also to be the scientifically cultivated man, then his
claim
% --- p. 30 ---
\opage{30}to be regarded as a figure cast in one piece must outrage every
truth-loving mind.

It stands written, after all, that death came into the world through sin. The
orthodox understanding of this is that death stands in no essential relation to
the human organism: had the first human beings not sinned, their bodies would
have been imperishable, indissoluble, unassailable by sickness and decay. But
they sinned, and with that this state of perfection was forfeited. Modern
science has long been clear not only that death struck manifold animal species
before human beings were formed, but that death stands in a necessary relation
to organic life. Here, then, science and orthodoxy flatly contradict one
another. Against orthodoxy's ``death came into the world with sin'' science
sets its ``no! death did not come into the world with sin'', just as against
orthodoxy's ``the world was created in 6 days'' it sets its ``no! the world was
not created in 6 days, since thousands of years passed before the earth became
what it is''\fnmark{*)}{When Prof.\ N.\ (Grundid.\ s.\ Logik II p.\ 456)
expresses himself as follows: It stands firm, provided only that one will not
let science's ``redemption'' juggle with faith, that pain and death existed for
millions of years before man appeared on earth. And it stands firm, provided
only that one will not let faith's ``redemption'' juggle with knowledge, that
death is the wages of sin and nature's corruptibility a transgression of
spirit'' --- then orthodoxy's hard proposition has here either evaporated into
a piece of pure wit, into a bold designation of the indisputable kinship
between sin and death, or else the reconciliation is wholly unintelligible. And
when Dr.\ Heegaard (Indledn.\ til Eth.\ p.\ 394) declares that one must hold on
to the 6 days, because one surely cannot be a believing geologist or a
geological believer, this is a poor and empty play on words; for whoever ---
be it as believer or as knower --- assumes anything whatever about the earth's
coming to be has an assumption that geology can correct.}. Now
% --- p. 31 ---
% The Danish note above is closed with an unpaired “ where the print never
% opened the quotation from Prof. N.; the English opens it, and balances.
\opage{31}Prof.\ Nielsen's claim is that one can without contradiction unite
two such propositions in one's consciousness. What is it, then, that brings
about this astonishing result? Answer: they are absolutely heterogeneous.

One scarcely knows what to wonder at most in this answer: \emph{at} such
propositions really being given out as absolutely heterogeneous, or \emph{at}
absolutely heterogeneous phenomena being able to be housed in one I at all, or
finally \emph{at} the absolute heterogeneity itself being the ground of the
peaceable relation between them. One grasps this much: that those propositions
are to be able to subsist side by side in amity because they are to be kept so
carefully apart that they never come into contact at all. And one sees,
perhaps --- in parenthesis --- a nemesis upon the dialectician who formerly
went to such extremes in his dialecticising, in the fact that here it is a
criticism that is to rescue the theory. One understands
% --- p. 32 ---
\opage{32}well enough, too, that the propositions' absolute heterogeneity may
be the ground of their not being able to contradict one another, since the
condition of their doing so would have to be that they could be compared, and
that presupposes homogeneity up to a certain degree. (Homogeneity of course
includes heterogeneity within it and is always relative; otherwise it would be
sameness.) But what one does not see at all is that absolute heterogeneity can
be the ground of two principles, utterances, propositions being reconciled. On
the contrary, it is a decisive ground for their not being reconcilable at all,
since they cannot so much as enter into relation with one another. Absolute
heterogeneity is another expression for absolute irreconcilability. Prof.\
Nielsen might therefore just as well have given his proposition this form:
absolutely irreconcilable principles can be reconciled, or can be not
absolutely irreconcilable --- all without contradiction. He has said it, but in
other words. When at the time, in ``Fædrelandet'', I branded it a
self-contradiction to have faith and knowledge united in the same consciousness
once they are determined as absolutely heterogeneous principles, people rose
from several quarters, ready to show me with confident superiority that my
attack on Prof.\ Nielsen rested on a complete misunderstanding: where I put an
\emph{although}, the Professor put a \emph{because}, and so
forth\fnmark{*)}{see Høffding: ``Philosophie og Theologie'' p.\ 41. R.\ Schmidt:
``Om Selvmodsigelsen o.\ s.\ v.'' p.\ 17. Heegaard: Indledn.\ t.\ d.\ r.\
Ethik p.\ 389.}. There is a way of inferring that is just as dangerous and at
least just
% --- p. 33 ---
\opage{33}as common as the hothead's ``I cannot understand that man,
consequently he is wrong'': it is the echoer's ``I understand this man,
consequently he is right''. The pupil begins by standing astonished, uncertain,
full of questions and doubts before his teacher's discourse; but when it then
dawns on him --- often all at once --- how the teacher actually conceives the
matter, in what way he wants what he has set out to be understood, then, in his
joy at grasping what was a riddle to many and so lately to himself, he becomes
at that very moment an adherent, and the more unconditionally the less
independent he is by nature. Such a pupil of course traces every objection to
his master's doctrine, in the most favourable case, to misunderstanding and
ignorance, and in the worst case to a moral depravity; and one must make some
allowance for his standpoint: that anyone might understand what the teacher
meant and still have something to object to it has not entered his boldest
dreams. So I am supposed not to have understood that the absolute heterogeneity
was precisely the condition of the much-discussed reconciliation. Although I
developed my view at that time in fair detail, and in particular did not ---
as Mr.\ Clemens Petersen has lately stated untruthfully --- withdraw when my
opponent ``made as if to carry the matter into a thoroughgoing discussion'', I
shall nevertheless take this point of dispute up again here. The matter is
quite simple. If two propositions such as
% --- p. 34 ---
\opage{34}a) ``the world was created in 6 days'' and b) ``the world was not
created in 6 days'' are absolutely heterogeneous, then this is the same as
saying that the one is an utterance of a being (essentia) that is absolutely
heterogeneous with the being of which the other is an utterance. Fiat
applicatio --- the R.\ Nielsen who states proposition a is absolutely another
than (heterogeneous with) the one who states proposition b; to call them both
by the same name is an innocent amusement, like Hostrup's conceit in
``Soldaterløier'' of calling all three heirs Peter. That Prof.\ Nielsen calls
himself by the same name in both cases is a caprice one may allow him, so long
as one takes care not to be led astray by it. Here, however, we have not merely
two persons, but two persons each belonging to a world absolutely heterogeneous
with the other's; for every splitting of a subject is a splitting of the whole
objective world that mirrors itself in the subject, since the subject is in the
last instance an utterance of that world. If the subject is divided, that comes
in principle of the whole world's having split (essentia) beforehand. One ought
not, then, to ask how matters stand with a man's self-consciousness when he
houses fundamental views that contradict one another in this way; for such a
man cannot exist. Absolute heterogeneity is absolute dualism; but absolute
dualism is a masquerade joke, a dummy concealing a real homogeneity and unity.

Since, from the standpoint once taken up, one does not go outside the features
that revelation comprises,
% --- p. 35 ---
\opage{35}one cannot exhibit its absurdity any further: the Christian
revelation does after all contain --- considering the new doctrine --- an
astonishing quantity of rational truths that will not be docked to fit the
paradoxes now so much in favour\fnmark{*)}{On this point the unhappiest abuse of
words has been carried on in our more recent literature, in that every truth
uttered in the New Testament is described as not having arisen in any human
being's heart or brain. Every proposition that human beings can grasp and
understand must, however, have been conceived by human beings in the first
place; for producing and grasping are one and the same thing seen from two
different sides. Next, every such truth is described as paradoxical, as opposed
to the human: thus it is supposed to be a paradoxical, a positively Christian
proposition that one should love one's enemies. But in reality it conflicts not
with understanding or reason at all, but with sensuousness; that one should not
repay evil with evil is a pure command of reason. If everything that comes into
conflict with our sensuous drive were paradoxical, then reason itself would be
the paradoxical.}. If, on the other hand, one places oneself --- as one is
entitled to do --- \emph{outside} the standpoint (Christian revelation) but
\emph{on} its fundamental proposition in that proposition's general form (the
possibility of absolutely heterogeneous utterances being united), there is
infinite play for the wildest attempts to satisfy the practical interest, the
personal need. Such attempts have been made already: witness the manifold
positive religions, which outdo actual fairy tales in fabulousness, and yet
precisely
% --- p. 36 ---
\opage{36}are the fruit of the deepest personal need. All of these Prof.\
Nielsen must let stand; whether they have a rational content of ideas or not
makes no odds from a supra-rational standpoint. One cannot even say: the madder
the better --- for everything may claim equal rank where personal need, purged
or cleansed of all reason, is the measure. That a man, as worshipper of his God
and devotee of his science, asserts absolutely conflicting statements is, from
the middle of the nineteenth century onward, no obstacle to his uttering the
clear, indeed the highest, truth. It is Denmark that has the honour of this
discovery, beside which one may perhaps set N.\ F.\ S.\ Grundtvig's, but
assuredly neither Ole Rømer's nor H.\ C.\ Ørsted's.

But one need not stop here; the positive religions have no exclusive right to
the satisfaction of the anti-rational need. It is my need and nobody else's
that I have to satisfy; and if I thereby come to live in a swarm of absolutely
heterogeneous utterances, what of it? It conflicts with reason, with thought.
Reason! Thought! They have no vote where the highest is at stake. One may
choose whichever contradictory assumption one likes; once one declares it
demanded by one's personal, one's ethical need, Prof.\ Nielsen can raise no
objection to it. To make this vivid, one single supra-rational instance.
Suppose a man has the personal faith that Martin Luther was an incarnation ---
not, let us say, of the
% --- p. 37 ---
\opage{37}whole world-principle, but of the Sirius-principle, say; and suppose
that for the same man, as knower, Luther was the well-known reformer who has
nothing whatever to do with the Sirius-principle. What good will it do then to
call such a thing unreasonable and unnatural? Probability and reason have
nothing to say on this point. Or would the philosopher really look for
reassurance in the thought that nobody will hit upon sheer absurdities? Would
he be reduced to appealing to probability --- that hereditary enemy of the
paradoxical --- to the probability that nobody will hit upon anything that
conflicts with reason to that degree? For reason is, after all, the only thing
that could have any objection to such assumptions; supra-reason would behold
everything, and behold, it was very good. Is the champion of principles really
to be reduced to wanting in principle, to the last extreme, what he most firmly
refuses to have in reality?

Since the whole strange doctrine is so hard to make plain when one goes
straight at the matter, adherents of Prof.\ Nielsen have from time to time done
the public the kindness of illuminating the relation by images or analogies.
Without dwelling on how precarious such illumination is by its very nature, we
shall simply see with what success it has been attempted. So as to leave
nothing relevant unmentioned, I will recall that Mr.\ Clemens Petersen some
years ago presented the relation between representation and concept in the
religious
% --- p. 38 ---
\opage{38}domain by the following image. When, it ran roughly, a man sees a
mill-sail turning on its axis, and one asks him whether the outermost or the
innermost part of the sail has the faster motion, then, if he goes by the
sensory impression, he will involuntarily answer: they go equally fast [?];
whereas grasped conceptually the outer end of course travels much faster than
the inner. Just as little, it went on, as one ought to ``torment representation
into seeing the concept'', just as little ought one to force representation's
illusion upon thought. And so too in religion: there is one truth for
representation, another, contradicting the first, for thought, and neither does
the other any damage. --- That was the parable. One sees that it answers to
Prof.\ Nielsen's theory, and that Mr.\ Petersen accordingly has --- as not many
know --- some other merit in this affair than that of being the Professor's
panegyrist. If one may draw an inference from his parable, the religious
representations turn out to be so many deceptions of the senses, of a doubtful
unavoidability. Mr.\ Petersen, however, seems to have meant to establish
something quite different.

Next, Dr.\ Heegaard (Indledn.\ til Ethiken p.\ 389) has supplied an
illuminating parallel. His words are these: ``At this (the contemplative) stage
in consciousness no conflict \emph{can} arise between faith and knowledge,
precisely \emph{because} they are heterogeneous, and not, as by a strange
misunderstanding has repeatedly (sic!) been maintained, \emph{although} they
are heterogeneous
% --- p. 39 ---
\opage{39}principles; for this `although' implies that the spheres
of faith and knowledge should be \emph{able} to bound one another, which is
impossible so long as consciousness hovers above them. Let us take an example:
the contemplative consciousness can without contradiction house at one and the
same time an aesthetic and an astronomical problem, not although but precisely
because they are heterogeneous; it cannot of course have them both in
consciousness at the same time, but it can immerse itself in either of them as
much as it likes without the one problem ever suffering damage from the other,
just because they point back to heterogeneous starting points.'' It should be
needless to add any commentary on this passage; it must be plain even to the
most easily beguiled that there is not so much as a trace of an analogy here.
One will notice, incidentally, that Dr.\ Heegaard carefully refrains from
calling the heterogeneity absolute. When I first contested Prof.\ Nielsen's use
of that expression, Messrs.\ Schmidt and Høffding vied with one another in
explaining to me that this was exactly where the secret lay: had the
heterogeneity been relative, it would have been quite another matter, and the
reconciliation would have been impossible. Dr.\ Heegaard, though, seems to have
discovered that by laying the stress on this point one passed out of Scylla
into Charybdis.

The third and most discussed parallel is the one Mr.\ Rud.\ Schmidt has drawn
between this relation and the relation between the truth of poetry and the
truth of reality. This analogy has rather more interest, since
% --- p. 40 ---
\opage{40}Kierkegaard has perhaps hinted at it in passing --- see his use of
Gorgias's saying that the deceived man is wiser than the undeceived, cf.\
Afsluttende Efterskrift p.\ 349. The question, then, is first and foremost
whether the two terms are absolutely heterogeneous. One must answer no at once;
for the denial lies already in the very expressions used. If both have truth,
then both are moments in the realm of truth, and so far is it from being
possible to speak here of absolute heterogeneity that there must even be an
element in which poetry and reality are absolutely homogeneous. For if truth is
a concept that really can be applied to them --- or, to use a more intimate
expression, a determination of their essence --- then it must permeate them
both wholly and entirely, permeate every other determination of theirs. Both
are then utterances of a totality that embraces them. So much in general. But
it comes out no less clearly when one takes hold of the relation in a definite
example. Mr.\ R.\ Schmidt has used one which Prof.\ Nielsen, by not disowning
it, has acknowledged as his own. He cites the folk belief that when the meadow
steams on a summer morning it is the bog-wife brewing: just as in this case
there is for the educated man no contradiction between the poetic belief and
the natural-scientific cognition, so too there is no contradiction between the
heterogeneous assumptions of faith and knowledge. But in the first place the
imaginary meadow and the real
% --- p. 41 ---
\opage{41}meadow are not absolutely heterogeneous, and consequently neither are
the corresponding conceptions. Separate the bog from the bog-wife --- absolutely,
of course --- and what is the latter then for the imagination? Folk belief's
personification is nothing but reality's reflex in the imagination, which must
work incessantly with the real in order to bring forth something that is not
reality and yet is reality. To understand this one must use dialectic; but in
this matter dialectic has been given over to the gods once and for all. In the
second place there is no analogy here. Since the meadow that folk poetry has
may in general be described as an imaginary meadow, while the one on which the
natural-scientific conception is directed is the real meadow, the two are not
related to one and the same object. One might then just as well call it a
contradiction of the same kind as the one between faith and knowledge when I
say at one and the same time that \emph{I} own Lake Fuur (a painting by N.\
N.) and that \emph{the state} owns Lake Fuur (the real one). It is another
matter if one maintains that the real meadow's steam is an effect of the
bog-wife's brewing, and at the same time says that it is just as fully an
effect of physical causes: then we have an analogy, a contradiction of the kind
in question. For when, say, orthodoxy assumes a human being's elevation from
the earth and gradual ascent into the dogmatic heaven, and science denies this
miracle, first by virtue of the law of gravitation and second because, however
high one may be lifted into
% --- p. 42 ---
\opage{42}the air, one never gets out of this world's astronomical heaven and
into the transcendent one, then it is the very same point, the very same
ascension, of which both speak. --- Nor, then, is the analogy we are seeking to
be found here. But leave that aside. Even if one is content, in the analogy's
interest, to stress the ease with which in the example used two heterogeneous
conceptions glide past one another without mutual influence, one forgets, on
the one hand --- and this is very bad --- that the advance of our cognition of
reality makes it impossible to use certain expressions of fantasy that were
once serviceable enough\fnmark{*)}{See Ørsted: Aanden i Naturen II.\ p.\ 4.
``For an imagination that has made its own a living, present picture of the
world system, the expression the foundation of the earth is no better
than the foundation of a well-hung chandelier, but if possible even less
fitting''.}; and on the other hand one overlooks that the so-called gliding
past one another is in reality nothing of the sort, but would rather have to be
determined as the one conception's penetration into the other's lacunae or
pores, made possible by the fact that neither of them is quite complete. This
is why, for instance, the poet has the freer hand with historical material the
less known it is to his public; if the public must be presumed familiar with
every small feature of the material, he will necessarily be very much bound. It
would, however, be what Kierkegaard calls ``running off into parentheses'' to
enter on a closer discussion of this in
% --- p. 43 ---
\opage{43}connection with the question that matters. The decisive thing is and
remains that the poetic illusion, as a complete unicum, is no analogy at all.
It is, be it noted, only poetic when it momentarily alternates between being
and being cancelled; it is, be it noted, only poetic in so far as it is
conscious. The birds that peck at painted berries, or the lazzaroni in Naples's
little theatres who give away the hero by their gestures when the villain has
hidden, have not attained the poetic illusion. The strangest thing, however, is
that the defenders of this analogy would be in the worst straits of all if one
were to hold them to it in earnest. For anyone who means to ascribe the fullest
reality to some definite religious representation has acted against himself as
soon as he sets it beside the purely imaginary --- be the latter never so
beautiful and never so significant, both far more beautiful and far less foggy
than the bog-wife's brew. If, on the other hand, the bog-wife is transformed
into a real person, then we have one of those anti-rational truths before which
Prof.\ Nielsen must consistently bow; and then the game is in full swing with
the two truths that Prof.\ Nielsen champions in plain Danish but combats when
they appear as the Middle Ages' duæ veritates.

If we now look back at what has been objected here to the new doctrine, what
significance is left to the whole array of absolutely heterogeneous principles,
% --- p. 44 ---
\opage{44}utterances, truths --- or whatever one is to call these ill-placed
beings? Might it not be with them as with the villages that Potemkin showed his
empress as evidence of the district's flourishing condition? They were, in the
proper sense, merely set up, for they did not exist. Subjectively regarded,
there cannot possibly be absolutely heterogeneous objects for us: considered as
causes, they would have to call forth absolutely heterogeneous utterances of
consciousness in a subject that would consequently be wholly split. Can then,
objectively regarded, any utterance whatever that issues from a principle of
existence, or from a sum of objects forming a totality, be absolutely
heterogeneous with any other whatever? Think it, and you fall into
self-contradiction, since you posit a whole and deny what follows from it. To
be is to be comprised within an existence; when absolutely heterogeneous
objects \emph{are}, then they are not absolutely heterogeneous. Make the trial
with whatever objects, utterances, fantasies or what you will, and you will
find that even if they exclude one another in one respect, they nonetheless do
not do so in another, at any rate in so far as they can both fall within one
subject. Try: ``a yellow symphony'' --- that cannot be said, and yet tone and
colour are essentially homogeneous as wave phenomena; ``an equilateral
right-angled triangle'' --- there is no such thing, and yet both adjectives
designate possible determinations of the
% --- p. 45 ---
\opage{45}concept triangle\fnmark{*)}{Not even such pronounced opponents as the
poets Mr.\ Henrik Scharling and Mr.\ Rud.\ Schmidt are absolutely
heterogeneous. On the contrary, there are points at which an absolute
homogeneity obtains between them. Mr.\ Schmidt hits upon such a point with
great assurance when in his book ``Er Hr.\ Lic.\ Henrik Scharling en alvorlig
Mand?'' p.\ 80 he advances the conjecture that Mr.\ Scharling has now ``as
usual written us something bad''.}. Take what has no significance whatever: the
meaningless, the non-existent. Separate it absolutely from what is, and it
cannot even be uttered; even this, which would seem to have a privilege of
existing in separation from its opposite, is nothing but a reflex of that
opposite. But when heterogeneity in the absolute sense has no power even here,
how much less has it power there, where existence and reality are ascribed to
everything.

One would do me a great injustice in supposing that in putting forward purely
elementary truths like these I have for a moment imagined myself to be teaching
a thinker like Prof.\ Nielsen something he did not know quite as well himself,
something he has not himself said long before I had any inkling of it. I owe it
in large part to Prof.\ Nielsen himself that I am in a position to contradict
him. But it is not so with truth that one who has communicated it as one of its
organs can afterwards take his gift back at pleasure, when as a weapon even in
the weakest hand it would be strong enough to
% --- p. 46 ---
\opage{46}disturb a boldly designed theory raised on a fragile foundation ---
a theory that will hamper spiritual advance, that will bar the way to every
effort to bring about harmony between reason and religion, between the human
spirit's intellectual and its ethical-religious demands, and that thereby
disturbs all work toward the solution of one of the human race's great tasks
for the future.

\begin{center}
\rule{2cm}{0.4pt}\\[-4pt]
\rule{2cm}{0.4pt}
\end{center}

%% ============================================================
%%  3. THE ANTI-RATIONAL  (printed pp. 47--54)
%% ============================================================
\chapter*{3.\ The Anti-Rational.}
\addcontentsline{toc}{chapter}{3. The Anti-Rational}
\markboth{The Anti-Rational}{The Anti-Rational}
\label{ch:3}

% --- p. 47 ---
\opage{47}One can tell oneself what must follow when a thinker overlooks those
elementary truths and, for the sake of a personal need that becomes --- once
universal validity is ascribed to it --- an idiosyncrasy tormenting and
oppressing others, attempts to hold fast to what cannot be held. The
consequence has shown itself already. It is an anarchy of concepts, a confusion
of concepts that undermines the ground on which our science is built, and of
which the theories of Dr.\ Heegaard are surely the most desperate expression
that any literature has to show in that line. For what have we now? We have the
well-known metaphysics and the well-known logic with their concepts: freeborn,
children of nature and reason, living their life in a world in which they are
ordered as members and where their essence unfolds and comes into its own. And
then on the other side we have a world where they must do the lowest bondsman's
service --- the service, namely, in which the bondsman must deny his origin and
his essence. For these same concepts and powers are employed in the
anti-rational world in such a way that they may express nothing
% --- p. 48 ---
\opage{48}rational at all. For if they do express something rational, they at
once transform the impotent powers for which they are to do forced labour ---
the so-called categories: the paradoxical, the anti-rational, the positive ---
into rational magnitudes, which is what those powers dread as certain beings in
the legend dread the daylight. Those children of light are therefore compelled
to put out their own light as well, so that the desired darkness may reign.
Take whichever concepts you like: principle, law, consequence --- they are
employed so that they shall express nothing of their essential content, indeed
so that they shall express \emph{nothing} at all, in order that they may
designate what their tyrants are. For the tyrants are nothing apart from those
they have subjugated, and even when they have subjugated them they are nothing
but wretchedness all the same: they must borrow everything from others --- the
shirt off their backs, as the saying goes --- and then they give themselves the
air of being the real capitalists of the spiritual life, or rather, they are
given out to be such.

The moment one really takes the doctrine of the anti-rational at its word, all
discussion stops, and indeed all communication. And this happens not because
one has failed to understand Prof.\ Nielsen properly, but precisely when one
has understood him properly: one could not misunderstand him more grossly than
by setting about to debate the matter with him, and he could not set a beginner
a harder test than by inviting him to a joint debate. That this is how things
stand I shall try to show --- taking care, of course,
% --- p. 49 ---
\opage{49}not to fall into the snare that the matter by its very nature lays
for anyone who gets involved with it.

In the first place, all communication is made impossible. When a freemason
wants to find out whether another man is a freemason, he does not ask him about
it, nor does he say that he is one himself; he uses no words at all to give and
receive this communication. He may talk of wind and weather, of trade and
politics, indeed he may even talk of freemasonry; but if he means to reach his
goal, he uses a means of communication lying quite outside language: he
contrives an opportunity to give the other's hand a peculiar pressure. If the
man in question is not a freemason, he notices merely a pressure, or notices
nothing at all; if he is a freemason, he understands the communication and
communicates in return. It stands similarly with the anti-rational religious
ethicist. He can no more communicate through language than the freemason can,
language being a system of signs for determinations of thought and reason; but
since on the other hand he cannot, like the freemason, use his hand instead of
the word, he chooses language's words to designate something that is as
independent of the rational determinations they designate as the freemason's
communication is of the handshake he gives. For that something exists only for
the organ of the anti-rational-religious, and does not exist at all for the
reason of which language is the expression and the means of communication. Now
if that something cannot
% --- p. 50 ---
\opage{50}be communicated through language, then of course it can as little be
communicated by any other means that stands under the law of the rational; the
path of the communication is a secret to everyone who does not possess the
organ. If one asks why language is used at all, this must doubtless --- for
that too is a secret --- be regarded as an incitement by which the secret
communication may take effect; if, on the other hand, one asks why language in
particular is used, the question is an unintelligent one, since the same
question could be raised about every other rational means, while about the
secret means no question can be raised at all, asking being a rational act.

Let us see the correctness of this in some examples. When, for instance, there
is talk of an anti-rational-religious ethics, then the word ethics, in the mouth
of the anti-rational-religious man and for his ear, does designate the rational
determination that the word expresses; but for the anti-rational-religious
organ it designates something quite different, which we --- precisely because
it falls outside all reason --- cannot possibly designate here. When it is said
further that the paradoxical ethics ``dialectically understood has its rational
side''\fnmark{*)}{Heegaard, op.\ cit.\ p.\ 439.}, then in the first place the
word dialectically cannot designate the rational determination of which it is
the expression; it designates a peculiarity of the anti-rational that can be
grasped only by the organ, and which the organ
% --- p. 51 ---
\opage{51}takes up into itself, occasioned by reason's apprehending what the
word expresses. In the second place the word rational can in no way carry its
accepted meaning; it merely indicates a further peculiarity, for whose reception
by the organ the determination ``the rational'' is an incitement. And when it
is said, finally,\fnmark{*)}{Heegaard p.\ 406, 439, 445.} that the anti-rational
in ethics emerges as a consequence of the rational, then the word consequence
here of course does not designate what it appears to designate, but yet another
new and interesting peculiarity of the anti-rational, to whose reception the
rational determination ``consequence'' incites. Indeed, this holds even of the
expression ``the anti-rational'' itself. It is anti-rational only for thought
and reason; so that for whose apprehension it is an incitement cannot be the
anti-rational, but is something about whose relation to thought, about whose
validity or invalidity for thought, there can be no question at all, since its
peculiarity is in each and every respect only for the organ. And so throughout:
the words with their rational determinations are so many incitements
occasioning the anti-rational-religious organ's many peculiar functions, of
which these words are in no way the expression; and whatever those functions
may be, communicated they cannot be.

The very thing that is the reason why the anti-rational cannot be communicated
also grounds its not being capable of objectification, and indeed of not being
thought
% --- p. 52 ---
\opage{52}at all, or made the object of any cognition in which thinking is a
moment. The reason is that it falls absolutely outside all rationality. That
the stress must fall on ``absolutely'' goes without saying, since nothing in
existence is exhausted by rationality, and since the rational everywhere
presupposes the real, without which it possesses no reality itself. But the
rational and the real permeate one another, so that they belong to one
another's essence; or, more precisely: it is one and the same essence that
distinguishes itself into its two essential sides. The rational therefore
stands under the law of the real, and the real under the law of the rational.
If, then, there shows itself a real that wholly withdraws from the law of the
rational, then of course no rational determination can express its essence in
the least degree, and it can exist neither for thinking --- which relates
itself to the ideal --- nor for the cognition that corresponds to the real,
because thinking is everywhere a moment in that cognition, just as the ideal is
everywhere a moment in the real. The determination ``the real'' can then be
applied to the anti-rational as little as any other; like those mentioned
earlier, it is only an incitement for the organ to function in some definite
direction. The same then holds of the word function, and of the remaining
expressions, all of which designate something in which there is a rational
element.

If the anti-rational is not understood in this way, it is transformed into a
rational; for a rational is a something that is real, a real is a something
that is
% --- p. 53 ---
\opage{53}rational, but an anti-rational is a something that is neither real
nor rational, that is not an anti-rational --- that may indeed be called so,
but is not really rightly so called, since it can have nothing whatever to do
with the expression by which it is named. From which it follows, too, that no
science can be formed about it. A system of thoughts can very well be formed,
but this will mean no more than a system of incitements to the anti-rational's
mysterious working, which those thoughts do not express at all --- to which
belongs the point that the word incitement likewise expresses nothing
concerning the relation between the rational and the anti-rational, since there
is no relation whatever. That no criticism can be given of the anti-rational's
presentation, since no presentation of it can be found; and that no discussion
can be conceived with the man who stands at the anti-rational's standpoint ---
of this one now needs only to be reminded, since criticism and discussion are
acts standing under the law of the rational.

If we recall, finally, that knowing is a rational act --- which is not the same
as a merely rational one --- then it follows of itself that we know no real
that is not at the same time rational. We therefore do not \emph{know} any
anti-rational, and know no organ for the anti-rational. Is it then not at all?
No; for being has a rational side, and to say that the anti-rational is, is to
say that it is rational.
% --- p. 54 ---
\opage{54}What, then, are we to say of it? We can say nothing; we can as little
affirm as deny anything about it. And one must remember above all that what has
been said of it here expresses nothing about its essence, and is absolutely and
unconditionally foreign to it.

\begin{center}
\rule{2cm}{0.4pt}\\[-4pt]
\rule{2cm}{0.4pt}
\end{center}

%% ============================================================
%%  4. THE CONSEQUENCES OF DUALISM  (printed pp. 55--76)
%% ============================================================
\chapter*{4.\ The Consequences of Dualism.}
\addcontentsline{toc}{chapter}{4. The Consequences of Dualism}
\markboth{The Consequences of Dualism}{The Consequences of Dualism}
\label{ch:4}

% --- p. 55 ---
\opage{55}When a small, almost invisible crack has appeared in a glass, it can
sometimes, after holding at the same point for a long while, run through all at
once and split the whole glass. When one of the human body's extremities is
attacked by a disease, the damage can eat its way from that starting point into
the body and attack its health at the root. So it has gone with the dualism
that showed itself early\fnmark{*)}{Phil.\ Propædeutik 1857 p.\ 77.} in Prof.\
Nielsen's view of the relation between faith and knowledge. For a long time it
seemed harmless to his theory, since it did not concern the theory itself; it
rested only on the fact that the theory gained no influence over his practice.
But now it has broken through so vigorously that it has penetrated into his
fundamental philosophical view itself\fnmark{**)}{Bemærkninger om theoretisk og
praktisk Erkjendelse. Nord.\ Universitetstidsskr.\ Tillægshefte til 10de
Aargang, 1866. Kjbh.}. In one sense this might seem strange: one would think it
could have no effect on
% --- p. 56 ---
\opage{56}the philosopher's doctrinal edifice that outside his science he
entertained an unscientific conviction. Even if Prof.\ Nielsen were to go so
far as to forbid us to apply the theory he has given us --- well, one might
say, let us at least be theoretical within theory. And if that were not the
intention, why else has the Professor so often and with such great pleasure
told us the anecdote of Laplace's answer to Napoleon's question? Would it not
be most remarkable if the philosopher who has so strongly stressed science's
protest against any dogma whatever should end, like Victor Cousin in France, by
offering his concept of God ``comme une base, qui peut porter la trinité
chrétienne''?\fnmark{*)}{se Grundid.\ s.\ Logik II.\ p.\ 343.} But the fact is
that precisely because the separation Prof.\ Nielsen means to install between
life and science is so chimerical, precisely for that reason the split in
existence, the ``psychological'' discord, reaches right into science, where it
ought not to be able to reach at all. In his ``Bemærkninger om theoretisk og
praktisk Erkjendelse'' the Professor has let the split pass over from the
religious domain to the ethical, and from the ethical to that of cognition.
Here it is then not merely religion and science, but acting and knowing, and
within cognition again the theoretical and the ``practical'' cognition, that
stand over against one another as absolutely heterogeneous. Matters have
% --- p. 57 ---
\opage{57}taken the same course here as with the paradox. About a paradox, an
absolute paradox in and for itself, one may think what one likes; but whatever
one thinks, it must surely be plain that a mediated paradox is an utter
monstrosity. One will remember, too, how ridiculous Kierkegaard found all
``glimpsing'' in the face of the absolute paradox (Afsl.\ Eftskr.\ p.\ 433);
but what is it other than the drollest ``glimpsing'' when we are
told\fnmark{*)}{Heegaards Fremstilling af Prof.\ Nielsens Lære, Indledn.\ til
Ethiken p.\ 439.} that ``the paradoxical standpoint emerges \emph{in virtue of
the same law of development as holds within the rational}'', and that ``rational
science doubles its point of view in such a way that it at once recognises the
paradox's existential incomprehensibility and, precisely in virtue of this
recognition, also discovers the law of development according to which science
gives a reflection of the paradox's genesis''*).
% `*)' above: the page's second, repeated reference to the same note.
Thus the relative paradox comes to prepare the absolute one and to pave the way
for it. While it still runs, in Grundid.\ s.\ Logik II.\ p.\ 278, that ``power
in its essence is just as rational as knowledge'', and that ``the infinite
\emph{reason} determines the logical \emph{by the anti-logical} and
conversely'', in Dr.\ Heegaard willing itself, in so far as it is something
incommensurable for theoretical knowledge (\emph{an anti-logical}), is called
``a rational \emph{paradox}'' (p.\ 441), and this fine concept
% --- p. 58 ---
\opage{58}thus forms the desired transition. As the paradoxical has here struck
roots right down into immanence's rational world, so too has absolute
heterogeneity.

It is, however, hard enough on this point to extract a sufficiently sharp and
decided statement from Prof.\ Nielsen. He does use strong words, but he is apt
to add a little subordinate clause or qualification that robs the words of
their ``resolute'' stamp again; he takes back with one hand what he has given
with the other. He fortifies the gulf between practice and theory; he even
seems to go so far as to want to deny that theory as such can be carried over
into practice at all, for --- as it runs --- ``what is known theoretically
cannot be known practically''\fnmark{*)}{Anf.\ Sted p.\ 69.}. This claim does
seem to be of an incautiousness that has at least this much in its favour, that
it is ``resolute''. But suppose someone now put the question to the Professor:
do you really say that practice cannot be derived from theory? Then he can
answer in his own words: by no means, I have only said that practice cannot
\emph{directly} be derived from theory\fnmark{**)}{P.\ 66.}. And who in all the
world would contradict that? In the same way the Professor declared in the
introduction to Grundid.\ s.\ Logik I that science cannot directly refute the
miracle. A valuable category,
% --- p. 59 ---
\opage{59}the category \emph{not directly} --- worthy of the man who so sharply
reproached speculative dogmatics with ``the spongy softness with which the
category degree has to figure everywhere, in every dilemma however decisive,
partly directly, partly indirectly''\fnmark{*)}{Phil.\ Propædeutik.
Efteraarshalvaar 1865 p.\ 173.}.

Further: the Professor blows the bridge between thought and will sky-high,
having first erased feeling's territory from his map and divided that weaker
power's possessions between those two autocrats. He not only ascribes to will
the same originality as to thought (p.\ 68), but lays such weight on the
refutation of the Hegelian error that will is merely a modification of thinking
that one would suppose absolute heterogeneity proved as soon as this
misunderstanding had been refuted. But the question here is not at all whether
will can be derived from thought or thought from will; it is whether both these
utterances of essence cannot be derived from one and the same fundamental
essence. Deny that, and we have dualism pure and simple. But here too the
decisive thing --- or rather the lack of any decision --- is that in reality it
is denied only in appearance. For if one objects that thought and will do after
all stand in unbroken relation to one another, the Professor merely answers:
why, I have never concealed that, ``but thought is not for that reason will,
nor will thought, and so forth'' (p.\ 67) --- which of course it could not
occur to any
% --- p. 60 ---
\opage{60}human being to deny. If one then, so as to get back to the centre,
tries turning the Professor's own words against him as follows: if ``religion
presupposes a revelation and revelation a miracle, while science denies the
miracle and rejects revelation'' (p.\ 63), then the form of the religious that
can be united with science can only be one without miracles and without
revelation --- then Prof.\ Nielsen answers bitterly that ``such expectorations
are more characteristic of the lively heads than of the thorough ones'' (p.\
73), notwithstanding that it is he himself who has set up the premises from
which this conclusion must necessarily be drawn. And if one then loses heart
and objects at the last only that the Professor's view must after all be an
open dualism, the author of the ``Bemærkninger'' can answer with justice: why,
I say so myself, I am pretty nearly a dualist --- for, as it runs, ``there are
in the life of spirit strong oppositions \emph{bordering} on dualism'' (p.\
66). Pretty-nearly-dualism is a relative of what Prof.\ N.\ calls
``middle-way theology'': one cannot catch it, being so supple that it slips out
of the hands of anyone who means to grasp it.

They are fine, true and inspired words that the author of Grundideernes Logik
utters in the second part, p.\ 457: ``A feeble idea naturally has feeble
extremes with a weak middle and a feeble conclusion, and that is why immanent
speculation is always so anxious about anything that might resemble a dualism.
The mighty idea, divine in
% --- p. 61 ---
\opage{61}essence, on the contrary has mighty extremes; its issuing in a
vigorous dualism is merely the presupposition of the infinite superiority with
which it completes the conclusion.'' But what does it help us that these words
are true, when in the psychological sphere the ``conclusion'' fails to arrive
and irreconcilable contradictions are set face to face with one another? We all
owe Prof.\ Nielsen thanks and recognition for the energy with which he has
pursued the ideal right out into the utmost extreme of its opposite, and
admiration for the brilliant confidence with which he has always felt assured
that it must after all prove possible to show the unity of even the most
conflicting things within the mighty idea. We hear him without resentment speak
scornfully of those who would take fright at the very shadow of a dualism ---
but only, it goes without saying, so long as the real dualism is not given out,
like the man in the fairy tale, for its own shadow. Once the teacher's
confidence has become over-confidence, the pupil's security gives way before
the feeling of a real sundering: a split that cannot be hidden under indefinite
turns of phrase like the ones just cited, a breach that can by no means be
riveted by the kind of dialectic that consists in forcing contradictory
expressions together, such as ``rational paradox'' or ``absolutely
heterogeneous utterances that turn out to be
relative''\fnmark{*)}{Bemærkningerne p.\ 65.}.

Fortunately the author of the ``Bemærkninger'' does not content himself
% --- p. 62 ---
\opage{62}with stating general propositions; he chooses definite examples to
illustrate what he has advanced. But, as was to be expected in so precarious a
matter, Prof.\ N.\ is strangely unlucky in the examples he gives. He expresses
himself (p.\ 69) in the following way: ``Categories such as freedom of choice,
choice, resolution and the like do indeed fall under thought's illumination, in
so far as they become clear in ethical cognition; but in principle they are
none the less categories of will, and their cognition is essentially practical.
What is known theoretically cannot be known practically; it is no use appealing
to the fact that the content known is the same and only the mode of cognition
different, for in every decisive cognition the content is determined by the
mode, and the mode by the content. That courage as a virtue can be known in
theory must be granted; --- but if the courage that is known theoretically
really were in each and every respect the same, the very same, as the courage
that is known practically, what would follow from that? It would follow that a
capable theoretical thinker, without venturing anything, without exposing his
life to any real danger, would have to be able to think himself into true
courage so thoroughly that he thereby at the same time thought true courage
into himself, and so, by objective thinking, really became courageous, however
cowardly he might otherwise be by nature. But that a man who is in reality a
coward can no more speculate his way to real courage, however theoretically he
may think, than a man without money can, though he
% --- p. 63 ---
\opage{63}were to think theoretically and objectively until he burst,
speculate his way to real money, stands in no need of proof.''

I shall not dwell further on this last example of the money that one is
supposed not to be able to speculate one's way to, though one thought until one
--- in Prof.\ N.'s forceful expression --- burst. I think one may maintain
without chicanery that any merchant, indeed any mere gambler on the exchange
who has speculated, could teach Prof.\ Nielsen that one can speculate oneself
into money, and out of it. Or is the point perhaps lodged in the words
``theoretically objectively''? I hardly think speculations on the exchange can
be reckoned to subjective, ethical-practical thinking. But what I mean to hold
to is the example of courage. For this example is a testimony beyond all others
to the want of clarity and sober reflection that characterises the whole new
doctrine. First a division is made of the categories into categories of will
and categories of knowledge; then it is said of the former that their cognition
is essentially practical; and then courage is named as an example. That Prof.\
Nielsen lays weight on this is clear from the fact that the example occurs not
only once in the ``Bemærkninger'' but, in the echo Dr.\ Heegaard, a full four
times, namely pp.\ 385, 399, 400, 405. Now if courage here meant the
intrepidity that can be inborn in one man and denied to another almost as a
physical disposition, then anyone would certainly be rightly
% --- p. 64 ---
\opage{64}held up to ridicule who believed that, however cowardly he might be
by nature, he could become courageous --- by nature --- through objective
thinking. But what was then said of courage as a \emph{practical} determination
would unfortunately hold in the same measure of all the theoretical endowments
too; for a man who is by nature slow, sluggish in thought and weakly gifted can
just as little speculate his way to inborn aptitude, acuteness or wit. But
since the express reference here, as we are told, is to courage as a
\emph{virtue}, the determination given is simply false. The theoretical
cognition of what truth enjoins, of what a relation of duty is, of what honour
requires, of what a fatherland means, can and should pass over into flesh and
blood in such a way that it combats and overcomes the natural timidity wherever
that stands in the way of doing what is right. Whoever denies this possibility
ought not to presume to plead life's cause against the theories; for he reveals
himself as an armchair philosopher, a pure theoretician, however much he may
extol and even deify the practical, \emph{however eagerly he strives to make it
the one thing needful.}

For rightly considered, what else is Prof.\ Nielsen doing? He speaks, to be
sure, of the double world, and often seems a pure dualist; but in the last
instance one will find that it goes with him as with all dualists: they are of
necessity dualists only in appearance, and when it comes to the point they are
driven to cast away one of the two opposed terms they have brought into
% --- p. 65 ---
\opage{65}unnatural tension with one another. It turns out that Prof.\ N.,
polemicising against making thought everything and thinking one with the whole
man (p.\ 68), in reality makes another of the human spirit's utterances ---
namely the life of the will --- into everything, into the whole; for if the
worst comes to the worst, that side may be sacrificed, but not this one (p.\
76). Hear the author's own words: ``He who sets thought upon the practical
cognition of a holy will may perhaps forfeit theory, but he wins life, spirit's
eternal life; for \emph{spirit's life is in the will}''. Is it a philosopher
speaking, or a priest, a thinker or a moralist? In a man who speaks with an
ethical or godly purpose in view one finds such an expression justified and
gladly approves it; but in a man who means to \emph{explain} existence to us as
a thinker?

One ought perhaps to wonder at it less, since quite corresponding aberrations
are met with in other literatures too. See for instance Taine: ``Les
philosophes francais du XIX\textsuperscript{e} siècle'' on Maine de Biran p.\
59: Il n'observe que le moi, l'âme, l'être, la substance. Où les trouve-t-il?
Dans la volonté\dots. Ainsi le moi n'est plus ce tout indivisible et continu,
dont nos idées, nos plaisirs, nos peines sont les parties composantes, isolées
par fiction et par analyse. c'est une force ou faculté, portion du tout, mise à
part, élevée au-dessus de toutes les autres. Le reste est mon bien: celle-ci
est moi''. In this way the witty Frenchman characterises this doctrine, which
makes
% --- p. 66 ---
\opage{66}the will into the whole man. But although the phenomenon is thus not
unique, I must confess that I for my part was highly astonished when, toward
the close of Prof.\ N.'s treatise, I came upon the passage that begins as
follows: ``Should it really still be impossible, as some maintain, for one and
the same consciousness to acknowledge both the practical and the theoretical,
faith and knowledge,'' --- it will be remembered that the Professor designates
by these words positive religion and the human spirit's intellectual demand ---
``then in any case it would be the theoretical, and not the practical, that
would have to be given up''. How so? Then the matter is not settled and proved
after all; it may still be possible that dualism cannot be carried through,
that one must give up saying both yes and no? And I who thought that this view
had --- in the author's opinion --- been utterly refuted as lacking the barest
rudiments of dialectic and psychology (see p.\ 65). This sudden concession of
the opposite's possibility has something very reassuring about it. The thing
now is to change tactics: it is no longer dualism that is to be attacked, but
the way out of it that it becomes our task to block. That way out is hinted at
in various expressions in Prof.\ Nielsen and his adherents; the manoeuvre
appears in various forms, but in reality it is easily recognised as one and the
same.

In this passage it is said first and foremost that if one of the two sides is
to be given up, it must be the
% --- p. 67 ---
\opage{67}theoretical. One notices that the ship is in distress and can be
saved only by throwing its ballast overboard: if all it thereby quite reliably
achieves is to capsize, its foundering is made tolerably evident. Regarded
quite generally, the way out chosen here has something very comic about it. The
whole doctrine consisted essentially in taking Kierkegaard's ``Øieblikket'' in
one hand and Ørsted's ``Aanden i Naturen'' in the other, clapping them together
and binding them into one. Now if someone comes forward with the objection that
this joining or tying together will not really do, and if he is stubborn enough
not to give way --- very well! then the other party will not be the quarrelsome
one, and cuts Ørsted away and keeps Kierkegaard. But what then remains of the
whole doctrine which has been described, so foolishly, as the greatest merit of
a man whose well-earned fame rests on quite other and far more genuinely
scientific achievements? What does it then contain that has not been said far
better, more richly and more fully by Kierkegaard?

If we now subject the matter to a closer examination: how does Prof.\ Nielsen
prove the possibility of giving up the theoretical? ``Existence,'' it runs,
``is practical; a consciousness can exist without being theoretical, but a
consciousness cannot be theoretical without existing. It is therefore possible
to put existence highest and to keep theory.'' Is this sophism really becoming
in a thinker, a thinker
% --- p. 68 ---
\opage{68}of Prof.\ Nielsen's rank? It is a play on words, a jesting toying
with an ambiguous term. First the word ``exist'' is taken in the purely
abstract sense of being at all, and then it is rightly said that a
consciousness cannot be theoretical without existing; then it is taken in the
peculiar sense just introduced, that of the life of the practical ideas, and
now it is said of existence that it must of course be put highest. Alas, this
``highest and lowest'' is a precious designation to Prof.\ Nielsen: if the
mutually contradictory utterances really could be reconciled, then it would
surely be one and the same thing on which shelf I put them. If faith's
superiority may not express itself in a reshaping of knowledge, it is not easy
to grasp what it consists in. But so it stands: because one cannot theorise
without existing, i.e.\ eating, drinking, sleeping etc., therefore the
paradoxes of a positive faith must be set above the rational results of the
human spirit. But never mind the weak grounding. Let us simply examine in and
for itself whether the theoretical really can be given up. Now, if theory
really were one with systematised science, and this in turn one with so and so
many printed leaves bound in paper or leather, then it would be feasible to let
science lie and let theory be theory. But if that is not how matters stand ---
if in the simplest conversation between a peasant and a fisherman, if in the
plainest utterance of two simple thoughts that hang together and in whose
% --- p. 69 ---
\opage{69}sequence there is consistency --- if, I say, in these quite
straightforward and everyday phenomena there is a germinating science, an
elementary theory, what then? No: science's essence depends neither on paper
and printer's ink nor on system and completeness. Theory is present wherever a
rational cognition takes place, be it ever so abrupt, ever so simple. Giving up
theory, then, is not so easy a business.

Another form of the same train of thought occurs in Dr.\ Heegaard, op.\ cit.\
p.\ 447. The presentation there is designed to characterise the ``will's
concentration out beyond the rational'' that is found in the believer in the
paradox. It is so strong that under it all regard for theory disappears; ``yes,''
we are told, ``not even rational ethics at its maximum can come into
consideration here''. Good! So long as the individual does not have this
concentration, he is only half a believer, and a half faith is no faith; but as
soon as he has raised himself to faith's height, he has lost all theory
entirely from view. What plainer statement of dualism's impossibility could one
wish for than this, given by its own founder and set down on paper by one of
its blindest defenders!

To this turn attaches another --- imposing in Kierkegaard, but unfortunate
here --- that the true believer has no time to occupy himself with
% --- p. 70 ---
\opage{70}scientific theory\fnmark{*)}{See Heegaard p.\ 395, smlgn.\ Nielsen:
Evangelietroen og Theologien p.\ 155.}. If this turn clears up nothing else, it
shows at least that our theoreticians, like the old viking kings who, content
with the honour, would readily give up their share of the booty, magnanimously
renounce for their own persons the faith which, as generals, they so bravely
win for everyone else.

I recall hearing Prof.\ Nielsen hint at yet another way out in his lectures;
since I have no printed document to rely on, I touch on it of course with a
certain reserve, though I do not doubt that Prof.\ N.\ will acknowledge having
followed the train of thought that matters here. The Professor discussed the
biblical passage where it is told that heaven opened and the spirit descended
like a dove --- the same passage that Mr.\ A.\ C.\ Larsen has cited in his book
``Samvittighed og Videnskab'' pag 59. The question was raised whether science
might here bring its profane criticism to bear, and call the believer to
account for how he manages to hold fast to a vault when he knows there is none,
to believe that it opened, to take the dove for anything more than a pure
symbol when he knows that the spirit cannot possibly shrink into a dove. That
was the question, though it was put in other words. The Professor then rebuked
the spiritlessness that
% --- p. 71 ---
\opage{71}clung to the single image, that failed to understand that for faith
these images were in flux and flow, in an incessant movement in which it was as
impossible for criticism to seize them as for the understanding to grasp them.
When one hears such a thing for the first time, one is ashamed of having been
so spiritless; one feels like a chicaner put to shame. But on closer
reflection it turns out here too that the standpoint can be maintained only
when its distinctive character is given up and the strong tension between
faith's and knowledge's ways of seeing is slackened. Just now we saw that one
helped oneself by throwing theory aside; here it seems, conversely, to be faith
that has to pay. For the clear historical reality that the believer must ascribe
to the holy matters of fact, if he is not to find himself to his astonishment
in agreement with the rationalist, seems here to give way to a mystical
half-reality, an ideal reality that has nothing to say where the whole point is
existence and fact.

But back to Prof.\ Nielsen's sacrifice, or possible sacrifice, of theory for
the benefit of practical life. What is it, then, that is to be gained by it?
One does not usually make even a merely possible sacrifice for nothing at all.

It is the will's originality and self-sufficiency that is to be secured; so far
the sacrifice is intelligible enough. What is not intelligible is that anything
much can be supposed to be won by it. For it must surely stand firm that the
will in and for itself is
% --- p. 72 ---
\opage{72}quite as immanent, quite as little supernatural an utterance of what
is human as thought is. It is, after all, a view belonging to the naivety of a
bygone age that certain of human life's and certain of earthly life's phenomena
or utterances are more supernatural than the rest. It was on account of such a
view that many an old believing Jew shook his head when the oil in the
Copenhagen synagogue's eternal lamp was replaced by gas. For gas is a
revolutionary element, highly natural, highly intellectual, and the light it
spreads is strong and sharp; oil, by contrast, is a mild, thick-flowing liquid
with an ancient title as the one right bringer of light, and is besides related
to unction and all else of that kind. Is thought now to be determined in the
same way as profane, while the will is given out as being akin to the positive?
Is the will then faith, or the organ by which faith is grasped? So it is
claimed, but the claim is utterly false. The will is not somehow half and half
a transcendent and revealed power; it is purely human, purely rational. If
faith's object is a miracle, then faith itself is just as fully a
miracle\fnmark{*)}{Se ``Philosophiske Smuler'' pag.\ 95.}, and the will is no
closer to the miracle than thought is. But on this point Prof.\ Nielsen makes
himself guilty of a series of sleights of hand; and he who has so stubbornly
combated the sophisms of speculative dogmatics puts forward a
% --- p. 73 ---
\opage{73}sophism perhaps less dazzling than theirs, but no less to be
rejected. It runs in the ``Bemærkninger'': ``Should anyone, however, be so
lively in the head and so light in his thoughts that it really makes him giddy
to decide what is to stand highest and what lowest, faith or knowledge, and out
of sheer liveliness confuses highest with lowest and lowest with highest, let
him try to simplify matters by throwing away faith, the ideal of will, the
principle of character, and living on pure theory: he will not escape the
doubling even so''. What a pace! What strides, as though in seven-league boots!
From existence to faith, from faith to the ideal of will and the principle of
character! I am perhaps arrogating no excessive honour to myself in assuming
that it is my humble self that is meant by the giddy lively one, and that the
passage cited was called forth by a remark of mine in ``Fædrelandet'' of 16
February 1866. But if so, where have I spoken of throwing away first faith in
general and then the ideal of will and the principle of character? As little as
it is I who recommend the cure of amputating the head to save the life, so
little do I wish any other mutilation of what is human, by which nothing is
achieved but that the fever --- in Holberg's phrase --- leaves the patient.
What else have I
% --- p. 74 ---
\opage{74}tried to establish in my newspaper articles than that the union of
two absolutely heterogeneous principles in the same consciousness is
impossible? Where does it stand written that whoever gives up positive religion
dissolves the religious itself into the scientific? A personality such as the
American minister Theodore Parker is testimony enough to the looseness of that
claim. And how will anyone convince us that the man who breaks with orthodoxy
fails the ideal of will, that the man who wants to preserve character's unity
in the theoretical and in the practical life is throwing away the principle of
character? How will anyone establish that it is on this side that characterless
ness is to be found, and on that side firmness? Here, as so often, the new
doctrine must refer us to life and to what life teaches. What, then, does life
teach, what answer does it give? Why, here, as always, the difficulty is that
life teaches each man something different. But can anything be objected from
the new doctrine's side when a man declares for his own part that he has found
the truest characters and the purest will on the side where, according to the
theory, the ``intellectual amphibians'' were to be found, and on the other
hand the fluid spirits, vibrating with intellectual enthusiasm, precisely where
according to the theory character's firmness was to be found? Of his own accord
nobody would use so weak, so womanish an argument; but as far as I can see it
is Prof.\ Nielsen himself who has licensed its use.

There frequently occur in the world of spirit phenomena of so decisive an
importance that they have a claim on our interest, no matter whether they crop
up at ever so great a distance from us or in the midst of our own
% --- p. 75 ---
\opage{75}fatherland. There are others to which we would pay no attention had
they arisen in China or in America or in Germany, but which we feel obliged to
acknowledge or to contest for the sole reason that their nearness in time and
space equips them with a power that their inner worth would never of itself
confer on them. Prof.\ Nielsen's doctrine of the relation between religion and
science belongs to the latter class. The reason we Danes ought to heed it is
that it is Danish. And yet --- what am I saying, Danish! Is this unnatural
doctrine, at odds with itself, Danish? It is so only in the literal sense of
the word; in the deeper and truer sense it is very far indeed from being
capable of being called Danish, so long as one credits Danish thinking with the
qualities of soundness and clarity, so long as one agrees that the Danish cast
of thought is not afraid of the dark, as the French so often is, but still less
obscure, as the German so frequently is, and least of all shuns the light or
defies the law. It is Danish, then, only in respect of its origin, and in so far
as its coming into being is intelligible only from given historical
presuppositions on Danish ground. This relation to the past has been invoked in
order to assert the new doctrine's necessity, in order to present every
opposition to it as an impotent attempt to hold back the age's own mighty
wheel; but, busy as one was with feeling the pulse of world history, one has
forgotten that the question here is not
% --- p. 76 ---
\opage{76}whether the theory is comprehensible as a natural offspring of an
earlier condition, or explicable as an escape from a double fire, but whether
it is satisfactory, whether it is true, whether the pomp with which its
adherents proclaim it is in any way justified. And one has failed to consider
that every eclectic doctrine can with the same ease establish its genuine
descent and its hereditary right to truth's throne.

\begin{center}
\rule{0.16\textwidth}{0.4pt}
\end{center}

\emph{Postscript:} A piece by Prof.\ R.\ Nielsen that appeared after this book
had gone to press, ``Svar til Hr.\ Dr.\ phil.\ Pastor Zeuthen'', gives me no
occasion to add anything to what I have written or to correct it. For the most
part it contains only a reprinting of the author's sallies against Martensen's
Dogmatics.

\end{document}
